\documentclass{article}
\usepackage{polyglossia}
\setdefaultlanguage{arabic}
\setotherlanguage{english}
\usepackage{amsmath}
\usepackage{fontspec}
\setmainfont{Amiri}

\title{نظرية الأعداد الجبرية وتطبيقاتها}
\author{أحمد محمد}
\date{٢٠٢٦}

\begin{document}
\maketitle

\section{مقدمة}

تعد نظرية الأعداد الجبرية من أهم فروع الرياضيات البحتة. في هذا البحث ندرس
حلقات الأعداد الصحيحة الجبرية وخصائص التحليل الأولي فيها.

% RTL-001: missing bidi package — this document should load bidi for proper RTL
% The polyglossia package handles basic RTL but explicit bidi is recommended

% RTL-002: RTL/LTR boundary without proper markers
نستخدم نظرية Galois لدراسة التوسعات الحقلية. يعتبر حقل الأعداد الجبرية
\begin{english}$\mathbb{Q}(\sqrt{-5})$\end{english}
مثالاً كلاسيكياً على حلقة لا تملك خاصية التحليل الأولي الوحيد.

% Incorrect: LTR content without proper environment markers
نظرية Dedekind تعطي حلاً لهذه المشكلة عبر المثاليات. خوارزمية Euclidean
تنطبق فقط على الحلقات Euclidean.

% Correct: with proper language switching
نظرية \textenglish{Dedekind} تعطي حلاً لهذه المشكلة عبر المثاليات.
خوارزمية \textenglish{Euclidean} تنطبق فقط على الحلقات الإقليدية.

\section{الأساسيات}

% LANG-010: language detection in mixed content
ليكن $K$ حقلاً عددياً بدرجة $[K:\mathbb{Q}] = n$.
نعرف حلقة الأعداد الصحيحة الجبرية كالتالي:
\[
  \mathcal{O}_K = \{ \alpha \in K : f(\alpha) = 0 \text{ لبعض } f \in \mathbb{Z}[x] \text{ أحادي} \}.
\]

% AR-002: Arabic hyphenation issues
المثاليـــــــات الأولية في حلقة الأعداد الصحيحة تلعب دوراً محورياً في النظرية.
التحليل إلى عوامل أولية في الأعداد الصحيحة العادية لا يتحقق دائماً.

\section{النتيجة الرئيسية}

\begin{theorem}
كل مثالي غير صفري في $\mathcal{O}_K$ يتحلل بشكل وحيد إلى جداء مثاليات
أولية:
\[
  \mathfrak{a} = \mathfrak{p}_1^{e_1} \cdots \mathfrak{p}_r^{e_r}.
\]
\end{theorem}

هذا التحليل الوحيد للمثاليات يعوض غياب التحليل الأولي الوحيد للعناصر.

% Numbers in Arabic context
العدد $p = 2$ ينشطر في $\mathbb{Z}[i]$ كالتالي: $(2) = (1+i)^2$.
الأعداد الأولية ٢ و٣ و٥ و٧ تمثل الحالات الأساسية.

\section{الخلاصة}

قدمنا عرضاً شاملاً لنظرية التحليل الأولي للمثاليات في حلقات الأعداد
الصحيحة الجبرية. هذه النظرية تشكل الأساس لدراسات أعمق في نظرية الأعداد
التحليلية وحدسية فيرما الأخيرة.

\end{document}
